\section*{Chapitre \ref{ch:b2:solid-liquid-diagrams} --- Diagrammes binaires solide--liquide}

\begin{solution}{exo:b2:solid-liquid-diagrams:1}
À \qty{1300}{\degreeCelsius}, le \omterm{def:b2:solid-liquid-diagrams:liquidus}{liquidus} est à 0.541 et le \omterm{def:b2:solid-liquid-diagrams:liquidus}{solidus} à 0.609.
0.40~: liquide seul. 0.58~: liquide (0.541) et \omterm{def:b2:solid-liquid-diagrams:solid-solution}{solution solide} (0.609). 0.70~:
\omterm{def:b2:solid-liquid-diagrams:solid-solution}{solution solide} seule.
\end{solution}

\begin{solution}{exo:b2:solid-liquid-diagrams:2}
Fraction de solide $(0.58 - 0.541)/(0.609 - 0.541) = 0.57$~: \qty{1.15}{mol} de solide,
\qty{0.85}{mol} de liquide. Nickel~: $1.15 \times 0.609 = \qty{0.70}{mol}$ dans le
solide, $0.85 \times 0.541 = \qty{0.46}{mol}$ dans le liquide~; total \qty{1.16}{mol}
$= 2.00 \times 0.58$.
\end{solution}

\begin{solution}{exo:b2:solid-liquid-diagrams:3}
À pression fixée, $v = n + 1 - \varphi$ (aucune réaction). (a) $1 + 1 - 2 = 0$.
(b) $2 + 1 - 2 = 1$. (c) Liquide et deux \omterm{def:b2:solid-liquid-diagrams:solid-solution}{solutions solides}~: $2 + 1 - 3 = 0$.
\end{solution}

\begin{solution}{exo:b2:solid-liquid-diagrams:4}
Une décroissance régulière jusqu'à \qty{232}{\degreeCelsius}, un palier horizontal pendant que l'étain
se solidifie, puis de nouveau une décroissance régulière, plus lente près de la température ambiante. Sur le
palier, liquide et solide coexistent à pression fixée, \omterm{def:b2:equilibrium-shifts:variance}{variance} nulle~: la chaleur
évacuée est fournie par la cristallisation, et la température ne peut pas varier
avant que le dernier liquide se soit solidifié.
\end{solution}

\begin{solution}{exo:b2:solid-liquid-diagrams:5}
Benzène~: $\ln x_B = -(9870/8.314)(1/269.6 - 1/278.64) = -0.1429$, $x_B =
0.867$. Naphtalène~: $\ln x_N = -(19\,100/8.314)(1/269.6 - 1/353.2) = -2.017$,
$x_N = 0.133$. La somme vaut 1.000~: \qty{269.6}{K} ($\qty{-3.5}{\degreeCelsius}$) est
la température eutectique, et le liquide eutectique contient 0.133 de naphtalène.
\end{solution}

\begin{solution}{exo:b2:solid-liquid-diagrams:6}
Le liquide refroidit jusqu'à la branche du \omterm{def:b2:solid-liquid-diagrams:liquidus}{liquidus} côté plomb, où apparaissent des cristaux de (Pb)~; le
liquide s'enrichit en étain, jusqu'à l'eutectique à \qty{182.2}{\degreeCelsius}
et 62.1~\%. Là, le liquide restant se solidifie en \omterm{def:b2:solid-liquid-diagrams:eutectic}{mélange eutectique}.
Fraction massique d'eutectique~: $(30 - 18.91)/(62.13 - 18.91) = 0.26$. En poursuivant
le refroidissement, la solubilité de l'étain dans (Pb) diminue (ligne en tirets)~: un peu de (Sn)
précipite à l'intérieur des cristaux de (Pb).
\end{solution}

\begin{solution}{exo:b2:solid-liquid-diagrams:7}
23.3~\% de NaCl en masse~: $23.3/76.7 = \qty{0.304}{kg}$ de sel par kilogramme
d'eau. À \qty{-25}{\degreeCelsius}, au-dessous de l'eutectique, aucun liquide ne peut exister,
quelles que soient les proportions~: le sel et la glace restent côte à côte à l'état solide.
\end{solution}

\begin{solution}{exo:b2:solid-liquid-diagrams:8}
$x_{\mathrm B} = 0.5$ se situe entre E$_1$ et le composé~: $\mathrm{AB_2}$
cristallise en premier, le liquide se déplace vers E$_1$ et s'y solidifie en
eutectique de A et de $\mathrm{AB_2}$~; à la fin, A + $\mathrm{AB_2}$.
$x_{\mathrm B} = 0.9$ se situe entre E$_2$ et B~: B cristallise en premier, le liquide
atteint E$_2$~; à la fin, $\mathrm{AB_2}$ + B.
\end{solution}

\begin{solution}{exo:b2:solid-liquid-diagrams:9}
Au bord arrière de la zone, le liquide se solidifie en un solide qui ne contient que 0.78
fois sa fraction molaire en cuivre~: le cuivre rejeté reste dans le liquide, qui
avance et s'enrichit~; le cuivre est ainsi balayé vers l'extrémité du barreau.
Avec un rapport proche de 1, chaque passage n'élimine qu'une fraction de l'impureté~:
il faut de nombreux passages, chacun partant du barreau laissé par le précédent.
\end{solution}

\begin{solution}{exo:b2:solid-liquid-diagrams:10}
Démonstration comme dans le chapitre. Pour un liquide dilué, $\ln x_A \approx -x_B$, $1/T -
1/T^* \approx -\Delta T/T^{*2}$, donc $x_B = \Delta_{\text{fus}}H_A\Delta T/(RT^{*2})$~;
avec $x_B \approx n_B/n_A = bM_A$, $\Delta T = (RT^{*2}M_A/\Delta_{\text{fus}}H_A)\,b$.
Benzène~: $K_f = 8.314 \times 278.64^2 \times 0.07811/9870 = \qty{5.11}{K.kg/mol}$.
\end{solution}

\begin{solution}{exo:b2:solid-liquid-diagrams:11}
Pour \qty{100}{g}~: $19.0/24.31 = \qty{0.782}{mol}$ de Mg, $81.0/207.2 =
\qty{0.391}{mol}$ de Pb~; rapport $2.00$~: $\ce{Mg2Pb}$.
\end{solution}

\begin{solution}{exo:b2:solid-liquid-diagrams:12}
Le palier est proportionnel à la masse de liquide eutectique~: l'alliage inconnu en contient
$120/300 = 0.40$. Côté plomb, $w = 18.91 + 0.40 \times (62.13 -
18.91) = 36.2~\%$ d'étain~; côté étain, $w = 96.59 - 0.40 \times (96.59 - 62.13) =
82.8~\%$. La première rupture de la \omterm{def:b2:solid-liquid-diagrams:cooling-curve}{courbe de refroidissement} (plus haute pour l'alliage à 36~\%,
dont le \omterm{def:b2:solid-liquid-diagrams:liquidus}{liquidus} est sur la branche raide du côté du plomb) ou une micrographie (cristaux primaires
de (Pb) ou de (Sn)) permet de les distinguer.
\end{solution}

\begin{solution}{pb:b2:solid-liquid-diagrams:1}
\textbf{1.} Plomb \qty{327.5}{\degreeCelsius}, étain \qty{231.9}{\degreeCelsius}.
\textbf{2.} Pour \qty{100}{g}~: $62.13/118.71 = \qty{0.5234}{mol}$ d'étain et
$37.87/207.2 = \qty{0.1828}{mol}$ de plomb~; $x_{\mathrm{Sn}} = 0.5234/0.7062 =
0.741$.
\textbf{3.} Points (0, 327.5), (100, 231.9), (62.13, 182.2), (18.91, 182.2),
(96.59, 182.2).
\textbf{4.} Comme sur la figure~: liquide~; L + (Pb)~; L + (Sn)~; (Pb)~; (Sn)~; (Pb) + (Sn).
\textbf{5.} L (62.1~\% de Sn) $\longrightarrow$ (Pb) (18.9~\%) + (Sn) (96.6~\%).
\textbf{6.} $v = 2 + 1 - 3 = 0$~: la température et les trois compositions sont
fixées.
\textbf{7.} La \omterm{def:b2:solid-liquid-diagrams:solid-solution}{solution solide} riche en plomb (Pb)~: le plomb dissout l'étain, si bien que le solide en
équilibre avec le liquide contient déjà de l'étain, dans la proportion lue sur le
\omterm{def:b2:solid-liquid-diagrams:liquidus}{solidus}.
\textbf{8.} Il suit le \omterm{def:b2:solid-liquid-diagrams:liquidus}{liquidus} vers l'eutectique en s'enrichissant en étain,
jusqu'à 62.13~\%.
\textbf{9.} Liquide à 62.13~\% et (Pb) à 18.91~\% d'étain.
\textbf{10.} (Pb)~: $(62.13 - 40)/(62.13 - 18.91) = 0.512$~; liquide 0.488.
\textbf{11.} (Pb) à 18.91~\% et (Sn) à 96.59~\%~; (Sn)~: $(40 - 18.91)/(96.59 -
18.91) = 0.271$~; (Pb)~: 0.729.
\textbf{12.} Constituant eutectique 48.8~\% (le liquide présent juste au-dessus)~;
(Pb) primaire 51.2~\%. Le (Pb) de la question 11 est la somme des cristaux
primaires et des lamelles de (Pb) de l'eutectique.
\textbf{13.} Une rupture de pente au \omterm{def:b2:solid-liquid-diagrams:liquidus}{liquidus}, une décroissance plus lente, puis un palier à
\qty{182.2}{\degreeCelsius} qui dure 0.488 fois celui de l'alliage eutectique.
\textbf{14.} $\ln x_{\mathrm{Sn}} = -(7030/8.314)(1/T - 1/505.1)$.
\textbf{15.} $\ln x_{\mathrm{Bi}} = -(11\,300/8.314)(1/T - 1/544.6)$.
\textbf{16.} $x_{\mathrm{Sn}}(T) + x_{\mathrm{Bi}}(T) = 1$.
\textbf{17.} À \qty{110}{\degreeCelsius}~: $0.587 + 0.349 = 0.936$~; à
\qty{120}{\degreeCelsius}~: $0.621 + 0.382 = 1.003$~; à
\qty{130}{\degreeCelsius}~: $0.655 + 0.417 = 1.072$. La somme atteint 1 vers
\qty{119.5}{\degreeCelsius}~: \qty{120}{\degreeCelsius} au degré près.
\textbf{18.} $x_{\mathrm{Bi}} = 0.38$~; en masse, $0.38 \times 208.98/(0.38 \times
208.98 + 0.62 \times 118.71) = 0.52$.
\textbf{19.} Le modèle donne \qty{120}{\degreeCelsius} et 52~\% de bismuth, le
diagramme évalué \qty{138.8}{\degreeCelsius} et 57~\%~: environ
\qty{19}{\degreeCelsius} de trop peu.
\textbf{20.} L'étain qui cristallise n'est pas pur, mais une solution contenant du
bismuth~: son \omterm{def:b2:chemical-potential:chemical-potential}{potentiel chimique} est plus bas que celui de l'étain pur, le solide est
plus stable, et il cristallise du liquide à une température plus élevée pour une
composition donnée du liquide. La branche de l'étain du \omterm{def:b2:solid-liquid-diagrams:liquidus}{liquidus} est relevée, et elle
rencontre la branche du bismuth plus haut (un liquide non idéal ajoute sa propre correction).
\textbf{21.} Elle fond et se solidifie à une seule température, la plus basse du système~:
pas de domaine pâteux pendant lequel un joint déplacé se fissurerait.
\textbf{22.} Le plomb est toxique, et les déchets électroniques finissent dans l'environnement.
\textbf{23.} Des composants sensibles à la chaleur peuvent être soudés à plus basse température~;
mais le joint se ramollit à plus basse température en service, et le bismuth le rend
fragile.
\textbf{24.} \qty{570}{g} de bismuth et \qty{430}{g} d'étain.
\textbf{25.} Le modèle du liquide idéal prévoit un eutectique à $\boldsymbol{T_E
\approx \qty{120}{\degreeCelsius}}$, contre \qty{138.8}{\degreeCelsius}
pour l'eutectique évalué.
\end{solution}
